
\documentclass{article}
\usepackage{colortbl}
\usepackage{makecell}
\usepackage{multirow}
\usepackage{supertabular}

\begin{document}

\newcounter{utterance}

\centering \large Interaction Transcript for game `cladder', experiment `full\_v1.5\_default', episode 5310 with Qwen3.5{-}27B.
\vspace{24pt}

{ \footnotesize  \setcounter{utterance}{1}
\setlength{\tabcolsep}{0pt}
\begin{supertabular}{c@{$\;$}|p{.15\linewidth}@{}p{.15\linewidth}p{.15\linewidth}p{.15\linewidth}p{.15\linewidth}p{.15\linewidth}}
    \# & \multicolumn{2}{c}{Player} && \multicolumn{2}{c}{Game Master} \\
    \hline

    \theutterance \stepcounter{utterance}  
    & & & \multicolumn{4}{p{0.6\linewidth}}{
        \cellcolor[rgb]{0.9,0.9,0.9}{
            \makecell[{{p{\linewidth}}}]{
                \texttt{\tiny{[P1$\langle$GM]}}
                \texttt{You are an expert in causal inference. The following question is not a typical commonsense query, but rather a meticulously designed question created by a professor specializing in causal inference, intended to assess the students' mastery of the course content.} \\
\texttt{Imagine a self{-}contained, hypothetical world with only the following conditions, and without any unmentioned factors or causal relationships: Liking spicy food has a direct effect on occupation and salary. Skill has a direct effect on occupation and salary. Occupation has a direct effect on salary. Skill is unobserved. Method 1: We look at how liking spicy food correlates with salary case by case according to occupation. Method 2: We look directly at how liking spicy food correlates with salary in general. To understand how liking spicy food affects salary, is it more correct to use the Method 1 than Method 2?} \\
\texttt{Start your answer with 'Yes' or 'No', followed by additional reasoning or evidence to support your explanation.} \\
            }
        }
    }
     \\ \\

    \theutterance \stepcounter{utterance}  
    & \multicolumn{4}{p{0.6\linewidth}}{
        \cellcolor[rgb]{0.561,0.737,0.561}{
            \makecell[{{p{\linewidth}}}]{
                \texttt{\tiny{[P1$\rangle$GM]}}
                \texttt{No} \\
\\ 
\texttt{In this hypothetical world, **Skill** acts as a **confounder** because it has a direct causal effect on both the treatment (Liking spicy food) and the outcome (Salary), as well as on the mediator (Occupation). Since Skill is unobserved, we cannot adjust for it directly.} \\
\\ 
\texttt{**Method 2** (looking at the general correlation) estimates the **Total Effect** of liking spicy food on salary. This total effect includes the direct path (Spicy \$\textbackslash{}to\$ Salary) and the indirect path through Occupation (Spicy \$\textbackslash{}to\$ Occupation \$\textbackslash{}to\$ Salary). While the unobserved Skill creates confounding bias in this estimate, it captures the full causal impact of the treatment on the outcome.} \\
\\ 
\texttt{**Method 1** (stratifying by Occupation) attempts to control for Occupation. However, in this causal structure, Occupation is a **mediator** (a variable on the causal pathway from Spicy to Salary). Conditioning on a mediator blocks the indirect causal path (Spicy \$\textbackslash{}to\$ Occupation \$\textbackslash{}to\$ Salary). Consequently, Method 1 estimates only the **Direct Effect** of liking spicy food on salary, excluding the portion of the effect that operates through changing one's occupation.} \\
\\ 
\texttt{Furthermore, because Skill is an unobserved common cause of both Spicy and Occupation, conditioning on Occupation (a collider or a descendant of a collider in the context of the unobserved confounder) can induce **collider bias** (or M{-}bias depending on the specific graph topology, but here it primarily acts as blocking the path and potentially opening backdoor paths if Skill affects Occupation and Spicy). Even without complex bias, the fundamental error is that Method 1 answers a different question ("What is the direct effect?") than the question asked ("How does liking spicy food affect salary?", which implies the total effect).} \\
\\ 
\texttt{Therefore, Method 1 is not "more correct" for understanding the overall effect; it is actually incorrect for estimating the total causal effect because it artificially removes the indirect effect through Occupation. Method 2, despite potential confounding bias from the unobserved Skill, is the only approach here that attempts to capture the total relationship, whereas Method 1 structurally eliminates a significant part of the causal mechanism. If the goal is to understand the *total* effect, Method 1 is fundamentally flawed due to over{-}adjustment (controlling for a mediator).} \\
            }
        }
    }
    & & \\ \\

    \theutterance \stepcounter{utterance}  
    & & & \multicolumn{2}{p{0.3\linewidth}}{
        \cellcolor[rgb]{0.9,0.9,0.9}{
            \makecell[{{p{\linewidth}}}]{
                \texttt{\tiny{[GM$|$GM]}}
                \texttt{No} \\
            }
        }
    }
    & & \\ \\

    \theutterance \stepcounter{utterance}  
    & & & \multicolumn{2}{p{0.3\linewidth}}{
        \cellcolor[rgb]{0.9,0.9,0.9}{
            \makecell[{{p{\linewidth}}}]{
                \texttt{\tiny{[GM$|$GM]}}
                \texttt{no} \\
            }
        }
    }
    & & \\ \\

    \theutterance \stepcounter{utterance}  
    & & & \multicolumn{2}{p{0.3\linewidth}}{
        \cellcolor[rgb]{0.9,0.9,0.9}{
            \makecell[{{p{\linewidth}}}]{
                \texttt{\tiny{[GM$|$GM]}}
                \texttt{game\_result = WIN} \\
            }
        }
    }
    & & \\ \\

\end{supertabular}
}

\end{document}
